% Glossaire des modèles (type=modeles) de la base drone : détecteurs YOLO et variantes drone,
% autres détecteurs, backbones, suivi, navigation, vision-langage, early-exit, opérateurs.
% Chiffres : \citep[§]{paper_id} s'ils viennent d'une fiche ou de benchmarks.csv (+ \pp si
% preprint), sinon \horscorpus. Les parents sont définis avant leurs enfants.

% ---------------------------------------------------------------- Famille YOLO
\newglossaryentry{yolo}{type=modeles, name={YOLO}, sort={YOLO},
  description={\emph{You Only Look Once} : famille de détecteurs en une étape, qui prédisent en
  une seule passe, sur une grille, des boîtes, un score d'objet et des scores de classe, puis
  filtrent les doublons par NMS. Depuis YOLOv3, l'architecture se découpe en backbone, neck
  (fusion multi-échelle) et tête de détection. C'est le détecteur dominant du corpus drone :
  du Tiny YOLO quantifié sur Zynq (Tincy YOLO) aux YOLOv5 à YOLO26 sur Jetson, FPGA ou
  GPU de bureau, et la base de presque toutes les variantes \enquote{petits objets} (têtes P2,
  Transformers, convolutions Ghost).}}

\newglossaryentry{yolov2}{type=modeles, name={YOLOv2}, sort={YOLO 02}, parent=yolo,
  description={Deuxième version : backbone Darknet-19, boîtes d'ancrage (\emph{anchors}),
  normalisation par lot, couche \emph{passthrough}, une seule échelle de sortie
  \horscorpus. Dans le corpus drone, c'est le détecteur visuel du système multi-capteurs de
  \citet{arxiv-2007.07396-real-time-drone-detection-tracking} (caméras visible 416$\times$416 et
  infrarouge 256$\times$256) : F1 de 78,49 en visible et 76,01 en infrarouge
  \citep[§VI]{arxiv-2007.07396-real-time-drone-detection-tracking}, sur GPU de portable. Tiny
  YOLOv2 est cité par la revue de \citet{2022-ahmed-survey-uav-computing-platforms-hardware}
  pour son efficacité sur FPGA.}}

\newglossaryentry{yolov3}{type=modeles, name={YOLOv3}, sort={YOLO 03}, parent=yolo,
  description={Troisième version : backbone Darknet-53 (blocs résiduels), prédiction à trois
  échelles par un neck de type FPN, trois ancres par échelle, classes en sigmoïdes
  indépendantes. Face à YOLOv2 : multi-échelle et backbone plus profond, donc meilleurs petits
  objets. Dans le corpus : référence de détection de drone sur GPU, avec 72,3 d'AP sur Det-Fly
  pour 32\,ms sur RTX 2080\,Ti \citep[Table~II]{2021-zheng-air-air-visual-detection-micro} et
  97,7 d'AP50 sur Anti-UAV à 36\,FPS (plateforme n/r)
  \citep[Table~1]{arxiv-2103.13933-unmanned-aerial-vehicle-visual-detection}\pp ; mesuré sur
  Jetson Nano, Jetson Xavier NX et Raspberry Pi avec NCS2 par
  \citet{2022-feng-benchmark-analysis-yolo-performance-edge} : 1,7\,FPS sur Nano à 416, 6,1 sur
  Xavier NX \citep[Table~3]{2022-feng-benchmark-analysis-yolo-performance-edge}.}}

\newglossaryentry{yolov3-tiny}{type=modeles, name={YOLOv3-tiny}, sort={YOLO 03 tiny},
  parent=yolov3,
  description={Version réduite de YOLOv3 : suite de convolutions 3$\times$3 et de max-pooling,
  deux échelles de sortie au lieu de trois. Sur Jetson Xavier NX elle atteint 41,1\,FPS à 416
  pour 13,5\,W, contre 6,8\,FPS sur Nano \citep[Table~3]{2022-feng-benchmark-analysis-yolo-performance-edge}.
  Sur FPGA, \citet{2025-cali-optimized-implementation-yolov3-tiny-real} la réduisent encore (une
  seule tête 13$\times$13, filtres divisés par 5, W4/A4 par FINN) : 17,68 de mAP@0,5 sur la seule
  classe \enquote{car} de VisDrone-DET, 164,79\,FPS et 2,55\,W de carte à 200\,MHz sur XC7Z020
  \citep[Tables~9 et 14]{2025-cali-optimized-implementation-yolov3-tiny-real}. SATAY la porte
  en W8A16 sur Virtex UltraScale \citep[Table~III]{2023-montgomerie-corcoran-satay-streaming-architecture-toolflow-accelerating}.}}

\newglossaryentry{slimyolov3}{type=modeles, name={SlimYOLOv3}, sort={YOLO 03 z Slim},
  parent=yolov3,
  description={YOLOv3-SPP3 élagué par \emph{network slimming} : régularisation L1 sur les
  facteurs d'échelle des BatchNorm puis suppression des canaux sous un seuil global et un
  seuil local. À 832$\times$832, la variante -SPP3-90 retire 85,9\,\% des FLOPs et 87,5\,\% des
  paramètres pour 42,6\,\% de temps d'inférence en moins sur GTX\,1080\,Ti seulement
  \citep[§5]{arxiv-1907.11093-slimyolov3-narrower-faster-better-real}\pp : l'élagage de canaux
  ne suit pas les FLOPs en latence. Son principe est repris par
  \citet{2021-liu-real-time-small-drones-detection} pour YOLOv4.}}

\newglossaryentry{yolov4}{type=modeles, name={YOLOv4}, sort={YOLO 04}, parent=yolo,
  description={Quatrième version : backbone CSPDarknet53, module SPP et neck PANet, tête de
  YOLOv3 à ancres, activation Mish, augmentations Mosaic. Dans le corpus, deux usages. Détection
  de petits drones : \citet{2021-liu-real-time-small-drones-detection} l'élaguent
  (canaux, couches, \emph{shortcut}) : 90,5 de mAP@0,5 à 69\,FPS contre 93,6 à 43\,FPS pour le
  modèle complet, sur Tesla P100 \citep[Table~4]{2021-liu-real-time-small-drones-detection}.
  Portage FPGA : \citet{2024-nguyen-fpga-soc-implementation-yolov4-flying} l'adaptent au DPU
  (Mish vers LeakyReLU, SPP 5/6/8) et l'élaguent à 0,6 sur ZCU104 : 82,41 de mAP@0,25 à
  29,0\,FPS contre 83,70 à 17,9\,FPS, et 27,9 contre 30,1\,W à la prise
  \citep[Tables~11 et 13]{2024-nguyen-fpga-soc-implementation-yolov4-flying}.}}

\newglossaryentry{yolov4-tiny}{type=modeles, name={YOLOv4-tiny}, sort={YOLO 04 tiny},
  parent=yolov4,
  description={Version réduite de YOLOv4 (backbone CSPDarknet53-tiny, deux têtes). Nguyen et al.
  y ajoutent une troisième tête pour les petits objets (variante \enquote{3L}) : 72,24 de
  mAP@0,25 à 125\,FPS sur ZCU104 (DPU, 4 threads), contre 83,70 à 17,9\,FPS pour YOLOv4 complet
  \citep[Table~11]{2024-nguyen-fpga-soc-implementation-yolov4-flying}. Sur Jetson Nano à 416,
  \citet{2022-feng-benchmark-analysis-yolo-performance-edge} mesurent 15,0\,FPS contre 1,6 pour
  YOLOv4 \citep[Table~6]{2022-feng-benchmark-analysis-yolo-performance-edge}.}}

\newglossaryentry{yolov5}{type=modeles, name={YOLOv5}, sort={YOLO 05}, parent=yolo,
  description={Version d'Ultralytics, sans article : réimplémentation PyTorch à ancres, backbone
  CSP, neck PANet, activation SiLU, tailles n à x par multiplicateurs de profondeur et de
  largeur \horscorpus. C'est la base la plus utilisée des YOLO drone du corpus : TPH-YOLOv5,
  SF-YOLOv5, GCL-YOLO, GLAD, GL-YOMO. \citet{2022-jung-improved-yolov5-efficient-object-detection}
  y remplacent SiLU par ELU (\enquote{ConvELU}) : 95,5 contre 94,6 de mAP@0,5 sur un jeu
  propre \citep[Tables~2--3]{2022-jung-improved-yolov5-efficient-object-detection}.
  \citet{2023-zheng-detection-localization-tracking-multiple-mavs} l'emploient pour exclure les
  objets courants (RTX 3070).}}

\newglossaryentry{yolov5n}{type=modeles, name={YOLOv5n}, sort={YOLO 05 1n}, parent=yolov5,
  description={Taille \emph{nano}. SATAY la déploie sur ZCU104 en 640$\times$640 : 26,7 de mAP50
  COCO, 21,41\,ms et 14,5\,W (méthode de mesure n/r), contre 32,28\,ms et 8,58\,W pour le
  Jetson TX2 \citep[Table~IV]{2023-montgomerie-corcoran-satay-streaming-architecture-toolflow-accelerating}.}}

\newglossaryentry{yolov5s}{type=modeles, name={YOLOv5s}, sort={YOLO 05 2s}, parent=yolov5,
  description={Taille \emph{small}. SATAY la couvre sur Virtex UltraScale en W8A16 (VCU110 :
  46,4\,ms, 24,8\,W) \citep[Table~III]{2023-montgomerie-corcoran-satay-streaming-architecture-toolflow-accelerating}.
  Elle sert de détecteur de base à GLAD et à la variante ConvELU de Jung et al.}}

\newglossaryentry{sf-yolov5}{type=modeles, name={SF-YOLOv5}, sort={YOLO 05 z SF},
  parent=yolov5,
  description={\emph{Small-Fast} : YOLOv5s sans la couche C5 ni la sortie P5, avec un neck
  PB-FPN (chemin supplémentaire haut vers bas) et un SPPF devant chaque tête
  \citep[§3]{2022-liu-sf-yolov5-lightweight-small-object}. À l'inverse des variantes à tête P2,
  l'allègement vient du retrait de la tête des grands objets : 2,24\,M paramètres, 13,8\,GFLOPs,
  34,3 de mAP@0,5 sur VisDrone, 11,5\,ms sur RTX 2060
  \citep[Table~4]{2022-liu-sf-yolov5-lightweight-small-object}. Aucun test embarqué.}}

\newglossaryentry{tph-yolov5}{type=modeles, name={TPH-YOLOv5}, sort={YOLO 05 z TPH},
  parent=yolov5,
  description={\emph{Transformer Prediction Head} : YOLOv5x à quatre têtes (ajout d'une tête
  P2 \enquote{tiny}), blocs Transformer (puis Swin) dans les têtes, module CBAM,
  multi-scale testing fusionné par WBF et ensemble de modèles. Au challenge VisDrone 2021, un
  ensemble atteint 39,18 d'AP \citep[§4.2]{arxiv-2108.11539-tph-yolov5-improved-yolov5-transformer}\pp,
  mais le modèle seul fait 7,36\,FPS (34,0 d'AP, VisDrone test-dev) sur RTX 3090 Ti
  \citep[Table~5]{2023-zhao-tph-yolov5-boosting-object-detection}. Référence de précision
  \enquote{YOLO drone SOTA}, jamais portée sur un accélérateur dans le corpus.}}

\newglossaryentry{tph-yolov5pp}{type=modeles, name={TPH-YOLOv5++}, sort={YOLO 05 z TPH2},
  parent=tph-yolov5,
  description={Extension de TPH-YOLOv5 : la tête P2 est supprimée et un module CA-Trans
  (\emph{cross-layer asymmetric transformer}, attention locale éparse limitée à un voisinage
  2$\times$2) transfère son information vers le niveau P3
  \citep[§3.3]{2023-zhao-tph-yolov5-boosting-object-detection}. Sur VisDrone test-dev : 33,1
  d'AP à 11,86\,FPS contre 34,0 à 7,36 ; sur UAVDT : 30,1 d'AP à 42,19\,FPS contre 26,9 à 25,12
  \citep[Tables~5--6]{2023-zhao-tph-yolov5-boosting-object-detection}.}}

\newglossaryentry{gcl-yolo}{type=modeles, name={GCL-YOLO}, sort={YOLO 05 z GCL},
  parent=yolov5,
  description={YOLOv5 allégé pour petits objets : backbone à convolutions Ghost (modules ELAG3,
  GSPPF), neck avec tête P2 et suppression de la tête P5, perte Focal-EIOU
  \citep[§3]{2023-cao-gcl-yolo-ghostconv-lightweight-yolo}. Quatre tailles N/S/M/L ; la version
  S fait 1,64\,M paramètres, 39,6 de mAP50 sur VisDrone à 53\,FPS (GPU de mesure n/r), contre
  32,7 pour YOLOv5-S \citep[Table~3]{2023-cao-gcl-yolo-ghostconv-lightweight-yolo}. Aucune
  mesure sur plateforme drone.}}

\newglossaryentry{glad}{type=modeles, name={GLAD}, sort={YOLO 05 z GLAD}, parent=yolov5s,
  description={\emph{Global-Local MAV Detector} : un YOLOv5s global sur l'image entière, un
  YOLOv5s local sur une fenêtre 300$\times$300 centrée par un filtre de Kalman, et un module de
  mouvement (homographie, différence d'images) en cas d'échec
  \citep[§III]{arxiv-2312.11008-global-local-mav-detection-challenging}\pp. Sur ARD-MAV
  (objets de 0,02\,\% de l'image en moyenne) : 80 d'AP@0,5 à 146,5\,FPS sur RTX 3070 et
  23,6\,FPS sur Jetson Xavier NX
  \citep[Tables~IV et VIII]{arxiv-2312.11008-global-local-mav-detection-challenging}\pp.}}

\newglossaryentry{gl-yomo}{type=modeles, name={GL-YOMO}, sort={YOLO 05 z GLYOMO},
  parent=yolov5,
  description={Suite de GLAD : YOLOv5 à quatre têtes P3--P6 (la tête ajoutée est P6, pas P2),
  C3Ghost, modules d'ASF-YOLO, plus détecteur de mouvement (Lucas-Kanade, template matching) et
  validation par filtre de Kalman, en pipeline global puis local
  \citep[§III]{arxiv-2411.02582-real-time-detection-small-uavs}\pp. Sur Drone-vs-Bird : 76,3
  d'AP (IoU 0,1) \citep[Table~II]{arxiv-2411.02582-real-time-detection-small-uavs}\pp, soit moins
  que RT-DETR (0,773) ; 21,6\,FPS sur Jetson Xavier NX
  \citep[§IV-E]{arxiv-2411.02582-real-time-detection-small-uavs}\pp.}}

\newglossaryentry{yolov7}{type=modeles, name={YOLOv7}, sort={YOLO 07}, parent=yolo,
  description={Septième version : blocs ELAN, re-paramétrisation (RepConv), têtes auxiliaires
  à étiquetage guidé \horscorpus. Base de deux variantes drone du corpus. Dans l'étude de
  \citet{2022-liu-sf-yolov5-lightweight-small-object}, YOLOv7 atteint 76,1 de mAP@0,5 en
  18,0\,ms sur RTX 2060 (jeu WIDER FACE) \citep[Table~3]{2022-liu-sf-yolov5-lightweight-small-object}.
  YOLOv7-tiny est le détecteur de Bo et al. cité par la revue d'\citet{arxiv-2504.11967-securing-skies-comprehensive-survey-anti}\pp.}}

\newglossaryentry{efficient-yolov7-drone}{type=modeles, name={Efficient YOLOv7-Drone},
  sort={YOLO 07 z Efficient}, parent=yolov7,
  description={YOLOv7 avec tête P2 ajoutée, tête P5 et stage 5 supprimés (moins 56,29\,\% de
  paramètres), convolution creuse à masque (TGM-CESC) et fusion descendante P2--P3--P4
  \citep[§3.2--3.4]{2023-fu-efficient-yolov7-drone-enhanced-object}. À 640$\times$640 sur VisDrone :
  32,1 de mAP@0,5:0,95 à 89,29\,FPS (RTX 3090) contre 28,6 à 212,76\,FPS pour YOLOv7
  \citep[Table~5]{2023-fu-efficient-yolov7-drone-enhanced-object} ; 39,9 à 1280
  \citep[Table~3]{2023-fu-efficient-yolov7-drone-enhanced-object}. \enquote{Efficient} vaut face aux
  méthodes coarse-to-fine, pas face à YOLOv7.}}

\newglossaryentry{f-yolov7}{type=modeles, name={F-YOLOv7}, sort={YOLO 07 z F}, parent=yolov7,
  description={YOLOv7 dont les blocs ELAN sont remplacés par un module \emph{Faster}
  (convolution partielle de FasterNet) et dont la perte CIoU devient FIoU (sigmoïde au lieu de
  l'arctan) \citep[§4]{2025-du-f-yolov7-fast-robust-real}. Sur DUT Anti-UAV : 92,80 d'AP@0,5 à
  169\,FPS (RTX 4090) ; avec soustraction de fond en prétraitement (Pre-F-YOLOv7, caméra
  supposée fixe) 98,80 à 159\,FPS \citep[Table~2]{2025-du-f-yolov7-fast-robust-real}. GPU de
  bureau seulement.}}

\newglossaryentry{yolov8}{type=modeles, name={YOLOv8}, sort={YOLO 08}, parent=yolo,
  description={Version d'Ultralytics : tête découplée et \emph{sans ancres}, blocs C2f
  (avec opération \emph{split}), SiLU, régression par Distribution Focal Loss, tailles n à x
  \horscorpus. Base des variantes drone Drone-YOLO, DEAL-YOLO, YOLO-GCOF et DroneScan-YOLO, et
  détecteur le plus comparé sur plateformes embarquées : Jetson Orin et Raspberry Pi 5 chez
  \citet{arxiv-2502.15737-performance-analysis-you-only-look} et
  \citet{arxiv-2409.16808-comprehensive-evaluation-deep-learning-object}\pp, DPU de ZCU104 chez
  Karki et al., SATAY sur Virtex.}}

\newglossaryentry{yolov8n}{type=modeles, name={YOLOv8n}, sort={YOLO 08 1n}, parent=yolov8,
  description={Taille \emph{nano}. Sur Jetson Orin NX (jeu propre de drones Tello), TensorRT INT8
  donne 65,83\,FPS à 71,9 de mAP50-95 contre 58,82\,FPS à 84,4 en FP16, pour 10,3 à 10,9\,W
  \citep[Table~4]{arxiv-2502.15737-performance-analysis-you-only-look} : la PTQ INT8 coûte 12,5
  points. Sur le DPU de ZCU104, la variante \enquote{DPU-aware} (SiLU vers ReLU, split vers
  convolutions 1$\times$1) perd 30 points sur KITTI après INT8 (0,64 vers 0,34 de mAP50-95)
  \citep[Table~II]{arxiv-2607.13106-no-attention-no-problem-dpu}\pp.}}

\newglossaryentry{yolov8s}{type=modeles, name={YOLOv8s}, sort={YOLO 08 2s}, parent=yolov8,
  description={Taille \emph{small}. Orin NX : 56,32\,FPS en INT8 contre 48,24 en FP16, 11,0 à
  12,8\,W \citep[Table~4]{arxiv-2502.15737-performance-analysis-you-only-look} ; Raspberry Pi 5
  (NCNN FP32) : 7,32\,FPS. SATAY la déploie sur Virtex UltraScale (VCU110 : 122,8\,ms)
  \citep[Table~III]{2023-montgomerie-corcoran-satay-streaming-architecture-toolflow-accelerating}.
  Base de DroneScan-YOLO.}}

\newglossaryentry{yolov8m}{type=modeles, name={YOLOv8m}, sort={YOLO 08 3m}, parent=yolov8,
  description={Taille \emph{medium} (44 de mAP COCO) évaluée avec n et s sur Jetson Nano,
  Orin Nano et Raspberry Pi 5 avec accélérateur Hailo-8
  \citep[§IV-C]{arxiv-2409.16808-comprehensive-evaluation-deep-learning-object}\pp.}}

\newglossaryentry{yolo-obb}{type=modeles, name={YOLO-OBB}, sort={YOLO 08 z OBB},
  parent=yolov8,
  description={Variantes à boîtes orientées (\emph{oriented bounding boxes}) de YOLOv8n,
  YOLO11n et YOLO26n, qui ajoutent un angle à la régression de boîte, pour les images
  aériennes. Karki et al. les compilent aussi pour le DPU de ZCU104 : 56 de mAP50-95 sur DIOR-R
  et 31 sur DOTA pour YOLOv8n-OBB après INT8
  \citep[Table~II]{arxiv-2607.13106-no-attention-no-problem-dpu}\pp ; YOLO26n-OBB tient
  29,55\,FPS de bout en bout \citep[§IV-C]{arxiv-2607.13106-no-attention-no-problem-dpu}\pp.}}

\newglossaryentry{drone-yolo}{type=modeles, name={Drone-YOLO}, sort={YOLO 08 z Drone},
  parent=yolov8,
  description={Famille de cinq modèles (nano à large) dérivés de YOLOv8 : tête P2 (PAFPN à quatre
  niveaux), fusion \enquote{sandwich} et sous-échantillonnage RepVGG
  \citep[§3]{2023-zhang-drone-yolo-efficient-neural-network}. La tête P2 apporte l'essentiel du
  gain (mAP50 de 43,7 à 50,5 sur VisDrone val) mais coûte cher : sur Jetson TX2, la version
  nano met 127,4\,ms contre 30,8\,ms pour YOLOv8n à nombre de paramètres égal
  \citep[Tables~6--7]{2023-zhang-drone-yolo-efficient-neural-network}. Large : 40,7 de mAP50 sur
  VisDrone test \citep[Table~4]{2023-zhang-drone-yolo-efficient-neural-network}.}}

\newglossaryentry{deal-yolo}{type=modeles, name={DEAL-YOLO}, sort={YOLO 08 z DEAL},
  parent=yolov8n,
  description={YOLOv8n (0,994\,M paramètres) pour la détection d'animaux par drone : pertes
  Wise-IoU et NWD, P5 remplacée par P2, module SSFF, optimiseur SOAP, inférence en deux passes
  sur les détections peu confiantes \citep[§2]{arxiv-2503.04698-deal-yolo-drone-efficient-animal}\pp.
  93,3 de mAP@0,5 sur WAID \citep[Table~2]{arxiv-2503.04698-deal-yolo-drone-efficient-animal}\pp.
  Plateforme d'inférence n/r.}}

\newglossaryentry{yolo-gcof}{type=modeles, name={YOLO-GCOF}, sort={YOLO 08 z GCOF},
  parent=yolov8n,
  description={YOLOv8n pour le contre-drone à basse altitude : bloc GI-DGCST à la place de C2f,
  neck CSA-HSFPN, tête partagée Detect-OGCH, perte Focaler-SIoU
  \citep[§III]{2025-yu-yolo-gcof-lightweight-low-altitude}. 1,2\,M paramètres, 4,1\,GFLOPs ;
  92,6 de mAP@0,5 sur A\&T et 91,5 sur DUT Anti-UAV, évalués sur GPU
  \citep[§IV-D7]{2025-yu-yolo-gcof-lightweight-low-altitude}. Pas de tête P2.}}

\newglossaryentry{dronescan-yolo}{type=modeles, name={DroneScan-YOLO}, sort={YOLO 08 z Scan},
  parent=yolov8s,
  description={YOLOv8s à 1280\,px (10,23\,M paramètres) avec élagage dynamique de filtres
  (RPA-Block, masque non structuré), branche P2 légère (MSFD) et perte SAL-NWD
  \citep[§3]{arxiv-2604.13278-dronescan-yolo-redundancy-aware-lightweight}\pp. 55,3 de mAP@0,5
  sur VisDrone à 96,7\,FPS sur GPU de portable
  \citep[Table~4]{arxiv-2604.13278-dronescan-yolo-redundancy-aware-lightweight}\pp ; les
  paramètres élagués ne sont pas retirés physiquement
  \citep[§5]{arxiv-2604.13278-dronescan-yolo-redundancy-aware-lightweight}\pp.}}

\newglossaryentry{yolov9}{type=modeles, name={YOLOv9}, sort={YOLO 09}, parent=yolo,
  description={Neuvième version : backbone GELAN et information de gradient programmable (PGI),
  dont la branche auxiliaire est retirée à l'inférence \horscorpus. Dans le corpus, seulement
  dans \citet{2024-gao-dataset-detection-method-unmanned-aerial} : DotD-YOLOv9-C, avec une perte
  fondée sur la \emph{Normalized Distance} (DotD), atteint 92,6 de mAP@0,5 sur le jeu Anti2 de
  l'article \citep[Table~5]{2024-gao-dataset-detection-method-unmanned-aerial}, sans mesure de
  vitesse.}}

\newglossaryentry{yolov11}{type=modeles, name={YOLO11}, sort={YOLO 11}, parent=yolo,
  description={Version d'Ultralytics qui succède à YOLOv8 : blocs C3k2 et module d'attention
  C2PSA dans le backbone, tête sans ancres \horscorpus. L'attention et SiLU ne passent pas sur la
  DPU : Karki et al. approchent C2PSA par des convolutions 1$\times$1 et un hard-sigmoid
  \citep[§III-B]{arxiv-2607.13106-no-attention-no-problem-dpu}\pp.}}

\newglossaryentry{yolov11n}{type=modeles, name={YOLOv11n}, sort={YOLO 11 1n}, parent=yolov11,
  description={Taille \emph{nano}. La variante DPU-aware perd peu après INT8 sur ZCU104 : 0,59 vers
  0,57 de mAP50-95 sur VOC, 0,60 vers 0,59 sur KITTI
  \citep[Table~II]{arxiv-2607.13106-no-attention-no-problem-dpu}\pp, nettement mieux que YOLOv8n.
  Base de YOLO-Drone.}}

\newglossaryentry{yolov11s}{type=modeles, name={YOLOv11-small}, sort={YOLO 11 2s},
  parent=yolov11,
  description={Taille \emph{small}. NOVA l'entraîne sur COCO avec un zoom adaptatif autour de la
  boîte précédente (entrée 320$\times$256) : 92\,Hz sous TensorRT FP16 sur Jetson Orin NX
  \citep[§IV-B]{arxiv-2506.18689-nova-navigation-object-centric-visual}\pp, détecteur du suivi de
  cible en vol.}}

\newglossaryentry{yolo-drone}{type=modeles, name={YOLO-Drone}, sort={YOLO 11 z Drone},
  parent=yolov11n,
  description={YOLOv11n dont la tête utilise des GhostConv et des blocs C2f à la place de C3k2
  (\emph{GhostHead}) ; le backbone est inchangé, pas de tête P2
  \citep[Fig.~4]{arxiv-2511.10905-yolo-drone-efficient-object-detection}. 30,4 de mAP@0,5 sur
  VisDrone, 1,8\,ms sur GPU Tesla T4 \citep[Table~2]{arxiv-2511.10905-yolo-drone-efficient-object-detection}.
  Aucun test embarqué.}}

\newglossaryentry{yolo26n}{type=modeles, name={YOLO26n}, sort={YOLO 26n}, parent=yolo,
  description={Taille \emph{nano} de YOLO26 (écrit aussi YOLOv26n), présentée comme détecteur
  de bout en bout sans NMS ni DFL \citep[§III-B3]{arxiv-2607.13106-no-attention-no-problem-dpu}\pp.
  Deux usages dans le corpus. Sur le drone intercepteur SCOUT, affiné sur 3\,931 images du drone
  cible puis couplé à ByteTrack : TensorRT FP16 sur Jetson Orin NX, 48\,FPS (20,8\,ms) contre
  22,5\,FPS sous PyTorch \citep[Table~IV]{arxiv-2609.21005-project-scout-interceptor-drone-perimeter}\pp,
  31,82\,FPS en vol avec 92,2\,\% de détections valides
  \citep[Table~III]{arxiv-2609.21005-project-scout-interceptor-drone-perimeter}\pp. Sur ZCU104,
  Karki et al. la compilent pour le DPU : 34,05\,FPS de bout en bout, le NMS étant conservé
  \citep[§IV-C]{arxiv-2607.13106-no-attention-no-problem-dpu}\pp.}}

\newglossaryentry{yolox}{type=modeles, name={YOLOX}, sort={YOLO X}, parent=yolo,
  description={Variante sans ancres, à tête découplée et affectation dynamique des positifs
  (SimOTA) \horscorpus. Dans le corpus, base de LCAM-YOLOX (détection de feux de forêt).}}

\newglossaryentry{lcam-yolox}{type=modeles, name={LCAM-YOLOX}, sort={YOLO X z LCAM},
  parent=yolox,
  description={YOLOX (neck FPN+PAN, trois têtes, 640$\times$640) avec attention de canal et
  spatiale par couche (LCAM, type CBAM) et Soft-NMS, pour la détection de feux de forêt par
  drone \citep[§III]{2025-elhence-fpga-accelerated-yolox-enhanced-attention}. Élagué à 30\,\% puis
  PTQ INT8 pour le DPU : sur Kria KV260, 195\,FPS et 9,72\,W (carte) à 78,11 de mAP, 387\,FPS à
  2 cœurs ; Jetson Nano 65\,FPS, STM32H7A3 15\,FPS (précisions n/r)
  \citep[§V-A, §V-D]{2025-elhence-fpga-accelerated-yolox-enhanced-attention}. Tableaux absents de
  l'extraction : chiffres lus dans le texte.}}

\newglossaryentry{tincy-yolo}{type=modeles, name={Tincy YOLO}, sort={YOLO z Tincy},
  parent=yolo,
  description={Tiny YOLO modifié par FINN-R : poids binaires et activations 3~bits (W1A3) sauf
  première et dernière couches (8~bits). 50,1 de mAP sur PASCAL VOC contre 57,1 pour le
  modèle de départ \citep[§5.1]{arxiv-1809.04570-finn-r-end-end-deep}. Premier YOLO quantifié sur
  SoC Zynq du corpus, rapporté en GOp/s : 133 sur PYNQ-Z1 (2,5\,W) et 288 sur Ultra96 (9,7\,W),
  sans FPS ni latence \citep[Table~5]{arxiv-1809.04570-finn-r-end-end-deep}.}}

\newglossaryentry{tinyissimoyolo}{type=modeles, name={TinyissimoYOLO}, sort={YOLO z Tinyissimo},
  parent=yolo,
  description={Détecteur pour microcontrôleur inspiré de YOLOv1 : entrée 88$\times$88,
  convolutions 3$\times$3 uniquement, grille de sortie 4$\times$4, INT8 par QAT, dimensionné
  pour les 442\,kB de poids du MAX78000. 58,5 de mAP sur un VOC réduit à 3 classes ; sur
  MAX78000, 5,5\,ms et 196\,µJ par inférence (MCU seul), contre 359\,ms sur STM32H7A3
  \citep[§III, Table~II]{arxiv-2306.00001-tinyissimoyolo-quantized-low-memory-footprint}. Seul
  détecteur de type YOLO sur MCU du corpus.}}

\newglossaryentry{skynet}{type=modeles, name={SkyNet}, sort={SkyNet},
  description={Réseau de détection mono-objet conçu de bas en haut pour DAC-SDC : un
  \enquote{Bundle} (convolutions depthwise 3$\times$3 et pointwise 1$\times$1, BN, ReLU6)
  empilé 6 fois, choisi par recherche sous contrainte de latence GPU et FPGA, avec bypass et
  tête YOLO à deux ancres \citep[§4--5]{arxiv-1909.09709-skynet-hardware-efficient-method-object}\pp.
  Vainqueur de DAC-SDC'19 : IoU 0,716, 25,05\,FPS et 7,26\,W sur Ultra96, 67,33\,FPS et
  13,5\,W sur TX2 \citep[Tables~5--6]{arxiv-1909.09709-skynet-hardware-efficient-method-object}\pp.
  Sert aussi de backbone à SiamRPN++ et SiamMask. Repris par plusieurs équipes DAC-SDC
  2020--2021, dont une version en streaming à 299,3\,FPS \citep[Table~5]{2024-zhang-comprehensive-analysis-dac-sdc-fpga}.}}

\newglossaryentry{ultranet}{type=modeles, name={UltraNet}, sort={UltraNet},
  description={Réseau de détection mono-objet de huit couches de convolution, quantifié en
  W4/A4, qui domine DAC-SDC à partir de 2020 avec SkyNet. Les équipes 2021--2022 le portent en
  architecture streaming avec \emph{DSP packing} : la version SEUer 2022 atteint 70,3 d'IoU,
  2\,020,6\,FPS et 1,41\,W sur ZU3EG \citep[Table~6]{2024-zhang-comprehensive-analysis-dac-sdc-fpga}.
  IoU mono-objet, non comparable au mAP multi-classe de VisDrone.}}

% ---------------------------------------------------------- Autres détecteurs
\newglossaryentry{ssd}{type=modeles, name={SSD}, sort={SSD},
  description={\emph{Single Shot MultiBox Detector} : détecteur en une étape à ancres
  (\emph{default boxes}) sur plusieurs cartes de résolutions décroissantes, suivi d'un NMS.
  Dans le corpus : TGIIF, un SSD élagué en 8~bits, gagne le volet FPGA de DAC-SDC'18 (IoU 0,62
  à 11,96\,FPS sur XC7Z020) \citep[§5.2.1]{2021-xu-dac-sdc-low-power-object} ; SSD512 est évalué
  sur MAV-VID (96,7 d'AP50 à 32,4\,FPS)
  \citep[Table~1]{arxiv-2103.13933-unmanned-aerial-vehicle-visual-detection}\pp et sur Det-Fly
  (78,7 d'AP en 84\,ms avec ResNet-50, 68,8 en 53\,ms avec MobileNetV2)
  \citep[Table~II, §V-B]{2021-zheng-air-air-visual-detection-micro}.}}

\newglossaryentry{ssd-mobilenetv2}{type=modeles, name={SSD-MobileNetV2}, sort={SSD 1 MbV2},
  parent=ssd,
  description={SSD à backbone MobileNetV2, décliné en trois largeurs (1,0$\times$, 0,75$\times$,
  0,5$\times$). Sur nano-drone (AI-deck GAP8, 160\,MHz, INT8), la version 1,0$\times$ tient
  1,6\,FPS pour 134,5\,mW à la carte \citep[Table~II]{arxiv-2301.12175-bio-inspired-autonomous-exploration-policies}\pp ;
  BlankSkip la mesure à 1,50\,FPS pour 55,5 de mAP (version 0,5$\times$ : 3,50\,FPS, 41,5)
  \citep[Table~2]{arxiv-2603.28149-blankskip-early-exit-object-detection}\pp. Détecteurs
  multi-objets du corpus sur GAP8.}}

\newglossaryentry{ssdlite-mobiledet}{type=modeles, name={SSDLite MobileDet},
  sort={SSD 2 MobileDet}, parent=ssd,
  description={Détecteur mobile à têtes SSDLite (convolutions séparables) et backbone MobileDet,
  recherché pour les accélérateurs mobiles \horscorpus. Alqahtani et al. mesurent 27 de mAP COCO
  contre 16 à 17 pour SSD MobileNet V1 sur Jetson Orin Nano, Jetson Nano et Raspberry Pi 5 avec
  Hailo-8 \citep[§IV-C]{arxiv-2409.16808-comprehensive-evaluation-deep-learning-object}\pp.}}

\newglossaryentry{retinanet}{type=modeles, name={RetinaNet}, sort={RetinaNet},
  description={Détecteur en une étape à ancres sur pyramide de caractéristiques, entraîné avec
  la \emph{focal loss} \horscorpus. Dans le corpus : 77,9 d'AP sur Det-Fly avec ResNet-50
  \citep[Table~II]{2021-zheng-air-air-visual-detection-micro}, 78,21\,\% de mAP en baseline de
  la détection de feu \citep[§V-A1]{2025-elhence-fpga-accelerated-yolox-enhanced-attention}, et
  base de QueryDet.}}

\newglossaryentry{querydet}{type=modeles, name={QueryDet}, sort={RetinaNet z QueryDet},
  parent=retinanet,
  description={RetinaNet à tête P2 dont les niveaux fins ne sont évalués qu'aux positions
  proposées par une carte de présence de petits objets (\emph{Cascade Sparse Query}, convolution
  éparse) \citep[§3.2]{arxiv-2103.09136-querydet-cascaded-sparse-query-accelerating}\pp. Sur
  VisDrone : 28,32 d'AP à 2,75\,FPS avec CSQ contre 28,35 à 1,16\,FPS sans
  \citep[Table~2]{arxiv-2103.09136-querydet-cascaded-sparse-query-accelerating}\pp. Il réduit le
  coût de la tête P2, au prix d'une convolution éparse peu adaptée au FPGA \deduit.}}

\newglossaryentry{faster-rcnn}{type=modeles, name={Faster R-CNN}, sort={Faster R-CNN},
  description={Détecteur en deux étapes : un réseau de propositions (RPN) génère des régions,
  puis une tête les classe et affine les boîtes \horscorpus. Plus précis mais plus lent que les
  détecteurs en une étape : 18,0\,FPS contre 32,4 pour SSD512 sur MAV-VID
  \citep[Table~1]{arxiv-2103.13933-unmanned-aerial-vehicle-visual-detection}\pp, 70,5 d'AP sur
  Det-Fly \citep[Table~II]{2021-zheng-air-air-visual-detection-micro}. Base de ClusDet.}}

\newglossaryentry{clusdet}{type=modeles, name={ClusDet}, sort={Faster R-CNN z ClusDet},
  parent=faster-rcnn,
  description={Faster R-CNN + FPN précédé d'un réseau de propositions de \emph{groupes} d'objets
  (CPNet), d'une régression d'échelle (ScaleNet) et d'une fusion itérative
  \citep[§3]{arxiv-1904.08008-clustered-object-detection-aerial-images}\pp : les groupes sont
  recadrés puis détectés à échelle normalisée. VisDrone val : 26,7 d'AP (ResNet-50, 273\,ms sur
  GTX 1080\,Ti) à 28,4 (ResNeXt-101, 773\,ms)
  \citep[Table~1]{arxiv-1904.08008-clustered-object-detection-aerial-images}\pp. Typique des
  méthodes de précision hors temps réel.}}

\newglossaryentry{ufpmp-det}{type=modeles, name={UFPMP-Det}, sort={UFPMP-Det},
  description={Détecteur de petits objets en deux temps : un détecteur grossier propose des
  régions, qui sont fusionnées et assemblées en une seule mosaïque (\emph{Unified Foreground
  Packing}), détectée en une passe par un réseau à tête multi-proxy (MP-Det)
  \citep[§3]{arxiv-2112.10415-ufpmp-det-accurate-efficient-object}\pp. VisDrone val : 36,6 d'AP
  (ResNet-50) à 39,2 (ResNeXt-101), 152\,ms sur GTX 1080\,Ti pour la configuration de la
  Table~3 \citep[Tables~1 et 3]{arxiv-2112.10415-ufpmp-det-accurate-efficient-object}\pp.}}

\newglossaryentry{sahi}{type=modeles, name={SAHI}, sort={SAHI},
  description={\emph{Slicing Aided Hyper Inference} : cadre indépendant du détecteur, qui
  découpe l'image en tuiles recouvrantes, détecte sur chacune puis fusionne par NMS
  \citep[§3]{arxiv-2202.06934-slicing-aided-hyper-inference-fine}\pp. Applicable sans
  réentraînement ; avec fine-tuning sur tuiles, FCOS atteint 38,5 d'AP50 sur VisDrone
  \citep[Table~1]{arxiv-2202.06934-slicing-aided-hyper-inference-fine}\pp. Multiplie le nombre
  d'inférences par image, ce qui l'éloigne du temps réel embarqué \deduit.}}

\newglossaryentry{detr}{type=modeles, name={DETR}, sort={DETR},
  description={\emph{DEtection TRansformer} : backbone CNN suivi d'un encodeur-décodeur
  transformer qui produit un ensemble fixe de boîtes à partir de requêtes apprises ; l'appariement
  hongrois supprime ancres et NMS \horscorpus. Dans le corpus, évalué sur GPU : 97,1 d'AP50 sur
  MAV-VID à 21,4\,FPS \citep[Table~1]{arxiv-2103.13933-unmanned-aerial-vehicle-visual-detection}\pp.
  RT-DETR (sa version temps réel) sert de comparaison à GL-YOMO sur Drone-vs-Bird (AP 0,773)
  \citep[Table~II]{arxiv-2411.02582-real-time-detection-small-uavs}\pp.}}

% ------------------------------------------------------------------ Backbones
\newglossaryentry{darknet}{type=modeles, name={Darknet}, sort={Darknet},
  description={À la fois le framework C d'entraînement de YOLO et la famille de backbones
  Darknet-19 (YOLOv2) et Darknet-53 (YOLOv3, résiduel). Dans le corpus : Darknet-53 est le
  backbone de YOLOv3 sur Det-Fly \citep[§IV]{2021-zheng-air-air-visual-detection-micro},
  le framework sert aux mesures sur Jetson sans TensorRT de
  \citet{2022-feng-benchmark-analysis-yolo-performance-edge}, et \citet{2024-nguyen-fpga-soc-implementation-yolov4-flying}
  convertissent leurs modèles Darknet vers Keras puis Vitis AI.}}

\newglossaryentry{resnet}{type=modeles, name={ResNet}, sort={ResNet},
  description={Famille de réseaux à connexions résiduelles (ResNet-18, -50, -101\ldots). Backbone
  par défaut des détecteurs de précision du corpus (ClusDet, QueryDet, UFPMP-Det, SSD512,
  Faster R-CNN, Dogfight) ; ResNet-18 allégé dans YOPOv2-Tracker
  \citep[§III-C]{arxiv-2505.06923-yopov2-tracker-end-end-agile}\pp ; le DroNet de Loquercio et
  Palossi en est un dérivé peu profond à connexions résiduelles
  \citep[§III-A]{arxiv-1805.01831-64mw-dnn-visual-navigation-engine}\pp. Tiny-PULP-Dronet en
  retire les by-pass \citep[§II]{arxiv-2407.02405-tiny-pulp-dronets-squeezing-neural}\pp.}}

\newglossaryentry{vgg16}{type=modeles, name={VGG-16}, sort={VGG16},
  description={Classifieur de 2014 à 13 convolutions 3$\times$3 et 3 couches denses, régulier
  mais lourd \horscorpus. Dans le corpus : réseau de contrôle de bout en bout du modèle de vol de
  \citet{2020-krishnan-sky-not-limit-visual-performance} (28\,Hz sur Jetson AGX Xavier, 10 sur
  TX2, 1,3 sur NCS \citep[§4]{2020-krishnan-sky-not-limit-visual-performance}), et backbone
  allégé du SSD vainqueur de DAC-SDC'18. Un SSD-VGG16 dépasse 80\,W sur FPGA selon la revue de
  \citet[§VI-A]{arxiv-2509.04153-real-time-fpga-cnns-detection}\pp.}}

\newglossaryentry{alexnet}{type=modeles, name={AlexNet}, sort={AlexNet},
  description={Classifieur de 2012 à 5 convolutions et 3 couches denses. Dans le corpus, backbone
  de tous les trackers siamois pour drone (SiamAPN++, HiFT, TCTrack), choisi pour la latence :
  3,4\,ms sur Jetson AGX Xavier contre 13,7 pour MobileNetV2, à cause du coût d'accès mémoire
  des réseaux mobiles \citep[§4.1, Table~1]{arxiv-2203.01885-tctrack-temporal-contexts-aerial-tracking}\pp.}}

\newglossaryentry{squeezenet}{type=modeles, name={SqueezeNet}, sort={SqueezeNet},
  description={Classifieur compact fondé sur des modules \enquote{fire} (compression
  1$\times$1, puis expansion 1$\times$1 et 3$\times$3). L'équipe SystemsETHZ de DAC-SDC le
  binarise (\emph{halfFire}) et lui ajoute une tête YOLO : IoU 0,49 à 25,97\,FPS sur XC7Z020
  en 2018 \citep[§5.2.1]{2021-xu-dac-sdc-low-power-object} ; en 2019 la variante RectHalfSqznet
  avec tête SSD atteint 0,553 d'IoU à 55,1\,FPS sur ZU3EG
  \citep[Table~3]{2024-zhang-comprehensive-analysis-dac-sdc-fpga}.}}

\newglossaryentry{mobilenet}{type=modeles, name={MobileNet}, sort={MobileNet},
  description={Famille de backbones légers fondés sur la convolution séparable en profondeur
  (depthwise puis pointwise 1$\times$1) ; V2 ajoute les blocs résiduels inversés à goulot
  linéaire \horscorpus. Présents dans le corpus de la détection embarquée (SSD) à la
  navigation (Loquercio).}}

\newglossaryentry{mobilenetv1}{type=modeles, name={MobileNetV1}, sort={MobileNet 1},
  parent=mobilenet,
  description={Première version, sans bloc résiduel. Backbone de SSD MobileNet V1 chez
  \citet{arxiv-2409.16808-comprehensive-evaluation-deep-learning-object}\pp : 16 de mAP COCO
  sur Jetson Orin Nano \citep[§IV-C]{arxiv-2409.16808-comprehensive-evaluation-deep-learning-object}\pp.
  L'équipe iSmart2 de DAC-SDC'18 l'associe à une tête YOLO (IoU 0,573, 7,35\,FPS sur XC7Z020)
  \citep[§5.2.1]{2021-xu-dac-sdc-low-power-object}.}}

\newglossaryentry{mobilenetv2}{type=modeles, name={MobileNetV2}, sort={MobileNet 2},
  parent=mobilenet,
  description={Blocs résiduels inversés à goulot linéaire. Dans le corpus : classifieur
  32$\times$32 drone / oiseau / fond après soustraction de fond
  \citep[§3.4]{2020-seidaliyeva-real-time-accurate-drone-detection}, backbone des SSD sur GAP8
  (Lamberti et al., BlankSkip), de RetinaNet allégé dans QueryDet (28,94 d'AP)
  \citep[Table~7]{arxiv-2103.09136-querydet-cascaded-sparse-query-accelerating}\pp et baseline
  de la détection de feu à 13\,FPS et 68,23\,\%
  \citep[§V-A1]{2025-elhence-fpga-accelerated-yolox-enhanced-attention}.}}

\newglossaryentry{mobilenetv3}{type=modeles, name={MobileNetV3}, sort={MobileNet 3},
  parent=mobilenet,
  description={Version issue d'une recherche d'architecture, avec modules d'attention
  squeeze-excitation \horscorpus. Extracteur de profondeur de la politique de vol à haute
  vitesse de \citet{arxiv-2110.05113-learning-high-speed-flight-wild}, suivi d'une convolution 1D
  et d'un MLP qui prédisent trois trajectoires : 24,3\,FPS (41,6\,ms) sur Jetson TX2 pour 60\,\%
  de succès en forêt réelle \citep[Table~1]{arxiv-2110.05113-learning-high-speed-flight-wild}.}}

% ------------------------------------------------------- Réseaux binaires (FINN)
\newglossaryentry{bnn}{type=modeles, name={BNN}, sort={BNN},
  description={\emph{Binarized Neural Network} : poids et activations sur 1~bit ($\pm 1$) ; le
  produit scalaire devient un XNOR et un comptage de bits, la BN suivie du signe un simple
  seuil, ce qui convient aux LUT d'un FPGA. Dans le corpus, c'est la base de FINN
  \citep{arxiv-1612.07119-finn-framework-fast-scalable-binarized} et FINN-R, dont descendent
  Tincy YOLO et, dans la même lignée de flot dataflow, le portage W4/A4 de YOLOv3-tiny de
  \citet{2025-cali-optimized-implementation-yolov3-tiny-real}.}}

\newglossaryentry{cnv}{type=modeles, name={CNV}, sort={CNV}, parent=bnn,
  description={Classifieur binaire de référence de FINN pour CIFAR-10 et SVHN : trois fois (deux
  convolutions 3$\times$3 et un max-pooling) avec 64, 128 et 256 canaux, puis deux couches
  denses. CNV-max atteint 21\,900 images/s à 80,1\,\% de précision sur CIFAR-10 sur ZC706
  (11,7\,W) ; CNV-fix, 11\,600 images/s pour une latence de 0,55\,ms
  \citep[Tables~3--4]{arxiv-1612.07119-finn-framework-fast-scalable-binarized}. Hors vision
  aérienne : sert de preuve de concept du flot FINN.}}

% ------------------------------------------------------------------- Suivi
\newglossaryentry{siamtracker}{type=modeles, name={Tracker siamois}, sort={Tracker siamois},
  description={Suivi d'objet unique par réseau jumeau : le gabarit (127$\times$127) et la zone de
  recherche (287$\times$287) passent par le même backbone, et une corrélation entre cartes donne
  la position de la cible \horscorpus. Famille dominante du suivi aérien du corpus, avec un
  backbone AlexNet pour tenir sur Jetson (SiamAPN++, HiFT, TCTrack) ; SiamRPN++ et SiamMask à
  backbone SkyNet sont évalués sur GOT-10k (36,4 et 39,0 d'AO)
  \citep[Tables~8--9]{arxiv-1909.09709-skynet-hardware-efficient-method-object}\pp.}}

\newglossaryentry{siamapnpp}{type=modeles, name={SiamAPN++}, sort={Tracker siamois 1 APN},
  parent=siamtracker,
  description={Tracker siamois à réseau de propositions d'ancres adaptatives à double
  caractéristique (APN-DF) et agrégation attentionnelle (AAN)
  \citep[§III]{arxiv-2106.08816-siamapn-siamese-attentional-aggregation-network}\pp. 56,0 d'AUC de
  succès sur UAV20L \citep[Fig.~4a]{arxiv-2106.08816-siamapn-siamese-attentional-aggregation-network}\pp ;
  34,9\,FPS sur Jetson AGX Xavier en vol réel
  \citep[§V]{arxiv-2106.08816-siamapn-siamese-attentional-aggregation-network}\pp.}}

\newglossaryentry{hift}{type=modeles, name={HiFT}, sort={Tracker siamois 2 HiFT},
  parent=siamtracker,
  description={\emph{Hierarchical Feature Transformer} : l'encodeur fusionne les cartes de
  corrélation de haute résolution (couches 3 et 4) par attention multi-têtes, le décodeur prend
  celle de la couche 5 à la place des requêtes d'objets
  \citep[§3.2]{arxiv-2108.00202-hift-hierarchical-feature-transformer-aerial}\pp, tête sans
  ancres. 78,7 de précision à 20\,px sur UAV123
  \citep[Table~1]{arxiv-2108.00202-hift-hierarchical-feature-transformer-aerial}\pp ; 31,2\,FPS
  sur Jetson AGX Xavier en test réel
  \citep[§5]{arxiv-2108.00202-hift-hierarchical-feature-transformer-aerial}\pp.}}

\newglossaryentry{tctrack}{type=modeles, name={TCTrack}, sort={Tracker siamois 3 TCTrack},
  parent=siamtracker,
  description={Tracker siamois qui injecte le contexte temporel : convolutions dont les poids sont
  recalibrés par les descripteurs des images passées (TAdaConv) et Transformer à mémoire
  constante (AT-Trans) \citep[§3]{arxiv-2203.01885-tctrack-temporal-contexts-aerial-tracking}\pp.
  58,2 d'AUC de succès sur UAVTrack112\_L à 125,6\,FPS sur PC (GPU n/r), 27\,FPS sur Jetson AGX
  Xavier en test réel \citep[§1, §5]{arxiv-2203.01885-tctrack-temporal-contexts-aerial-tracking}\pp.}}

\newglossaryentry{bytetrack}{type=modeles, name={ByteTrack}, sort={ByteTrack},
  description={Suivi multi-objets par association de détections, y compris celles de faible
  confiance, en deux passes \horscorpus. SCOUT l'associe à YOLO26n sur Jetson Orin NX (deux pools,
  seuil 0,50, confiance minimale 0,06) et ajoute un filtre géométrique de rejet des sauts de
  détection \citep[§III-B]{arxiv-2609.21005-project-scout-interceptor-drone-perimeter}\pp. Cité
  par la revue sur caméras événementielles comme suiveur de référence sur FRED
  \citep[§5]{arxiv-2508.04564-drone-detection-event-cameras}\pp.}}

\newglossaryentry{kalman}{type=modeles, name={Filtre de Kalman}, sort={Kalman},
  description={Estimateur récursif d'état à partir de mesures bruitées, qui alterne prédiction
  (modèle de mouvement) et mise à jour \horscorpus. Omniprésent dans le corpus comme brique de
  suivi : prédiction de la fenêtre locale de GLAD
  \citep[§III]{arxiv-2312.11008-global-local-mav-detection-challenging}\pp, validation des
  correspondances de GL-YOMO, correction de la dérive de la VIO dans Swift
  \citep[Methods]{2023-kaufmann-champion-level-drone-racing-deep}.}}

\newglossaryentry{ekf}{type=modeles, name={EKF}, sort={Kalman 1 EKF}, parent=kalman,
  description={Filtre de Kalman étendu : linéarise les modèles autour de l'estimation. Dans
  YOPOv2-Tracker, un EKF à vitesse constante assure la cohérence spatio-temporelle de la cible
  au sein du NMS des primitives de mouvement
  \citep[§III-B]{arxiv-2505.06923-yopov2-tracker-end-end-agile}\pp.}}

\newglossaryentry{ukf}{type=modeles, name={UKF}, sort={Kalman 2 UKF}, parent=kalman,
  description={Filtre de Kalman \emph{unscented} : propage des points sigma au travers du modèle
  non linéaire, sans le linéariser \horscorpus. NOVA en fait tourner un à 200\,Hz pour fusionner
  IMU et détections asynchrones \citep[§III-B3]{arxiv-2506.18689-nova-navigation-object-centric-visual}\pp ;
  NGPS, un UKF à 15 états fusionnant GPS visuel, VIO et IMU
  \citep[§III-E]{arxiv-2607.18936-ngps-gps-denied-aerial-geo}\pp.}}

% -------------------------------------------------------------- Navigation
\newglossaryentry{dronet}{type=modeles, name={DroNet}, sort={DroNet},
  description={CNN résiduel léger à deux sorties (angle de braquage en régression, probabilité de
  collision en classification), entraîné sur des images de voiture et de vélo pour la navigation
  d'un drone sans carte \horscorpus. Environ 41\,MMAC par image
  \citep[§III-A]{arxiv-1805.01831-64mw-dnn-visual-navigation-engine}\pp. C'est le réseau de
  navigation de référence des nano-drones du corpus (PULP-Dronet, Tiny-PULP-Dronet,
  roofline de Krishnan et al.).}}

\newglossaryentry{pulp-dronet}{type=modeles, name={PULP-Dronet}, sort={DroNet 1 PULP},
  parent=dronet,
  description={DroNet porté sur le SoC GAP8 : Fixed16 Q4.12 avec réentraînement QAT,
  max-pooling 2$\times$2, fine-tuning sur images de la caméra embarquée, tiling par AutoTiler sur
  8 cœurs \citep[§IV]{arxiv-1805.01831-64mw-dnn-visual-navigation-engine}\pp. 89,1\,\% de
  précision de collision ; 6\,FPS pour 64\,mW ou 18\,FPS pour 284\,mW
  \citep[§VI, Table~II]{arxiv-1805.01831-64mw-dnn-visual-navigation-engine}\pp. Gamme nano
  (mW).}}

\newglossaryentry{tiny-pulp-dronet}{type=modeles, name={Tiny-PULP-Dronet},
  sort={DroNet 2 Tiny}, parent=pulp-dronet,
  description={PULP-Dronet réduit : canaux multipliés par $\gamma=0{,}125$ et branches by-pass
  supprimées, car le réseau est creux (jusqu'à 92\,\% de neurones inactifs) et sur-apprend
  \citep[§II]{arxiv-2407.02405-tiny-pulp-dronets-squeezing-neural}\pp. 6,4\,kB et 1,5\,MMAC
  contre 320\,kB et 41\,MMAC ; 88\,\% de précision, 89\,FPS pour 34\,mW sur GAP8, contre
  10\,FPS à fréquence égale \citep[§III]{arxiv-2407.02405-tiny-pulp-dronets-squeezing-neural}\pp.}}

\newglossaryentry{pulp-frontnet}{type=modeles, name={PULP-Frontnet}, sort={Frontnet},
  description={CNN sans connexions résiduelles qui estime la pose d'une personne devant un
  nano-drone pour l'interaction humain-drone, quantifié en INT8 (NEMO) et déployé par DORY sur
  GAP8 \citep[§IV]{arxiv-2103.10873-fully-onboard-ai-powered-human}\pp. Trois variantes : la
  80$\times$32 tient 18,5\,FPS pour 8,6\,mW (puissance minimale) ou 134,7\,FPS pour 86,6\,mW
  \citep[§V-B]{arxiv-2103.10873-fully-onboard-ai-powered-human}\pp.}}

\newglossaryentry{nanoflownet}{type=modeles, name={NanoFlowNet}, sort={NanoFlowNet},
  description={Réseau de flux optique dense dérivé de STDC-Seg (convolutions séparables, filtres
  divisés par 4), guidé à l'entraînement par les frontières de mouvement, pour l'évitement
  d'obstacles d'un nano-drone \citep[§III]{arxiv-2209.06918-nanoflownet-real-time-dense-optical}\pp.
  5,57\,FPS en niveaux de gris 112$\times$160 sur GAP8 (INT8 PTQ) ; 141\,FPS et 7,122 px d'erreur
  de bout sur MPI Sintel sur GTX 1070
  \citep[Table~IV]{arxiv-2209.06918-nanoflownet-real-time-dense-optical}\pp.}}

\newglossaryentry{dreamerv3}{type=modeles, name={DreamerV3}, sort={DreamerV3},
  description={Apprentissage par renforcement avec modèle du monde : un RSSM prédit dynamique,
  récompense et continuation à partir d'images, et l'acteur-critique s'entraîne sur des
  trajectoires imaginées. Dream-to-Fly l'applique à la course de drones à partir de pixels
  64$\times$64 : seul DreamerV3 converge, en 20\,M d'interactions et environ 240\,h
  d'entraînement, là où PPO et SAC échouent
  \citep[§IV-A]{arxiv-2501.14377-dream-fly-model-reinforcement-learning}\pp. Validé en
  hardware-in-the-loop (calcul déporté), sans mesure embarquée.}}

\newglossaryentry{yopov2}{type=modeles, name={YOPOv2-Tracker}, sort={YOPOv2},
  description={Réseau unique de planification de trajectoire pour le suivi agile : entrée RGB-D
  160$\times$96, backbone ResNet-18 modifié, sortie de primitives de mouvement qui jouent le rôle
  d'ancres, entraîné par rétropropagation du coût de trajectoire sans démonstrations
  \citep[§III]{arxiv-2505.06923-yopov2-tracker-end-end-agile}\pp. Latence perception-action de
  8,2\,ms sur Jetson Orin NX en vol réel (3\,ms sur PC)
  \citep[§IV-B]{arxiv-2505.06923-yopov2-tracker-end-end-agile}\pp, suivi jusqu'à 6\,m/s.}}

\newglossaryentry{swift}{type=modeles, name={Swift}, sort={Swift},
  description={Système de course de drones qui bat des pilotes champions : détecteur de coins de
  portes (U-Net à 6 niveaux, 384$\times$384), VIO et filtre de Kalman pour l'état, et politique
  de contrôle MLP (2 couches de 128) entraînée par PPO en simulation
  \citep[Methods]{2023-kaufmann-champion-level-drone-racing-deep}. Sur Jetson TX2 : 40\,ms pour
  le détecteur (TensorRT FP16), 8\,ms pour la politique sur CPU
  \citep[Methods]{2023-kaufmann-champion-level-drone-racing-deep}. Référence de latence
  perception-contrôle en vol agile.}}

\newglossaryentry{superpoint}{type=modeles, name={SuperPoint}, sort={SuperPoint},
  description={Détecteur-descripteur de points d'intérêt appris (descripteurs de 256 dimensions)
  \horscorpus. Dans NGPS, il extrait les points de l'image du drone et de la tuile satellite avant
  l'appariement par LightGlue \citep[§III-C]{arxiv-2607.18936-ngps-gps-denied-aerial-geo}\pp.}}

\newglossaryentry{lightglue}{type=modeles, name={LightGlue}, sort={LightGlue},
  description={Réseau d'appariement de points d'intérêt à attention, qui peut s'arrêter tôt sur
  les paires faciles \horscorpus. Dans NGPS (RANSAC sur homographie, au moins 20 correspondances),
  il représente 82\,\% du temps GPU : le module tourne à 2,6\,Hz sous CUDA et 0,5\,Hz sur CPU,
  sur Jetson Orin NX \citep[§IV-D]{arxiv-2607.18936-ngps-gps-denied-aerial-geo}\pp. Le
  positionnement visuel atteint 2,94\,m de RMSE contre 10,37\,m pour la VIO seule
  \citep[§IV-E4]{arxiv-2607.18936-ngps-gps-denied-aerial-geo}\pp.}}

\newglossaryentry{vinsfusion}{type=modeles, name={VINS-Fusion}, sort={VINS-Fusion},
  description={Odométrie visuelle-inertielle par optimisation sur fenêtre glissante, monoculaire,
  stéréo ou avec GPS \horscorpus. Fournit l'état aux contrôleurs de YOPOv2-Tracker (quadrotor
  maison, Orin NX) \citep[§IV-A]{arxiv-2505.06923-yopov2-tracker-end-end-agile}\pp et la VIO
  monoculaire de NGPS (fenêtre de 10 images clés, 150 points KLT)
  \citep[§III-B]{arxiv-2607.18936-ngps-gps-denied-aerial-geo}\pp.}}

\newglossaryentry{depthanythingv2}{type=modeles, name={DepthAnythingV2}, sort={DepthAnythingV2},
  description={Modèle de fondation d'estimation de profondeur monoculaire à backbone ViT
  \horscorpus. NOVA utilise la variante ViT-S, alignée sur la disparité stéréo par un polynôme
  d'ordre 2 pour compléter la profondeur : 31,6\,ms par image sur Jetson Orin NX
  \citep[§IV-C]{arxiv-2506.18689-nova-navigation-object-centric-visual}\pp.}}

% ----------------------------------------------------------- Vision-langage
\newglossaryentry{clip}{type=modeles, name={CLIP}, sort={CLIP},
  description={Modèle image-texte entraîné par contraste, dont les embeddings relient une requête
  textuelle à une image \horscorpus. Dans SINGER, un champ sémantique CLIP est distillé dans la
  scène de Gaussian Splatting du simulateur
  \citep[§IV-B]{arxiv-2509.18610-singer-onboard-generalist-vision-language}\pp.}}

\newglossaryentry{clipseg}{type=modeles, name={CLIPSeg}, sort={CLIP 1 Seg}, parent=clip,
  description={Segmentation guidée par texte au-dessus de CLIP ViT-B/16 (86\,M de paramètres) :
  une requête et une image donnent une carte de chaleur. Étage de perception de SINGER sur Jetson
  Orin Nano : 3\,Hz sans optimisation, 12\,Hz en ONNX sur CUDA, qui borne la boucle de
  perception \citep[§V]{arxiv-2509.18610-singer-onboard-generalist-vision-language}\pp.}}

\newglossaryentry{singer}{type=modeles, name={SINGER}, sort={SINGER},
  description={Politique de navigation vers un but sémantique donné par texte, entièrement
  embarquée, entraînée par imitation d'un expert (RRT* + MPC) dans un simulateur photoréaliste
  \citep[§IV]{arxiv-2509.18610-singer-onboard-generalist-vision-language}\pp. En vol réel
  (30 essais, bureau factice) : but atteint à moins de 1\,m dans 76,67\,\% des cas avec le nord
  fourni, 66,67\,\% avec capteurs embarqués seuls, contre 53,33\,\% pour une baseline
  CLIPSeg + PD \citep[§VI-D]{arxiv-2509.18610-singer-onboard-generalist-vision-language}\pp.}}

\newglossaryentry{svnet}{type=modeles, name={SV-Net}, sort={SINGER z SV-Net}, parent=singer,
  description={Réseau visuomoteur léger de SOUS-VIDE réutilisé par SINGER : il prend en entrée
  les logits de CLIPSeg convertis en carte de chaleur, un historique de masse et de poussée, et
  sort poussée collective et vitesses angulaires à 20\,Hz
  \citep[§V]{arxiv-2509.18610-singer-onboard-generalist-vision-language}\pp. Taille en
  paramètres, latence et puissance propres : n/r.}}

\newglossaryentry{smolvlm}{type=modeles, name={SmolVLM-256M}, sort={SmolVLM},
  description={Modèle vision-langage compact : encodeur SigLIP, projection par \emph{pixel
  shuffle} et décodeur SmolLM2-135M \citep[§II]{arxiv-2609.25415-cloud-edge-split-profiling-onboard}\pp.
  Profilé par \citet{arxiv-2609.25415-cloud-edge-split-profiling-onboard}\pp sur un Jetson et un
  GPU cloud : 2,76\,s tout embarqué, 0,72\,s tout cloud, 2,09\,s avec coupure après la
  projection (moins 24\,\% face à l'embarqué), pour 15\,J à bord contre 26\,J
  \citep[§V-A1, Fig.~6 et 9]{arxiv-2609.25415-cloud-edge-split-profiling-onboard}\pp.}}

\newglossaryentry{litevla-h}{type=modeles, name={LiteVLA-H}, sort={LiteVLA-H},
  description={Modèle vision-langage-action de 256\,M de paramètres, à deux modes (action en 1 à
  2 tokens, sémantique en phrase) et ordonnanceur à deux cadences où l'action passe en priorité
  \citep[§3]{arxiv-2605.00884-litevla-h-dual-rate-vision}\pp. Sur Jetson AGX Orin en FP16 :
  50,65\,ms (19,74\,Hz) pour l'action, 149,90\,ms (6,67\,Hz) pour une phrase
  \citep[§6.1]{arxiv-2605.00884-litevla-h-dual-rate-vision}\pp, la latence étant dominée par le
  pre-fill multimodal. Surcoût du double régime : 3,6\,W.}}

\newglossaryentry{nerf}{type=modeles, name={NeRF}, sort={NeRF},
  description={\emph{Neural Radiance Field} : un MLP encode une scène 3D comme densité et couleur
  en tout point, rendue par intégration le long des rayons \horscorpus. Hors temps réel : dans le
  corpus, seulement porté à l'échelle d'un nano-drone par Tiny-DroNeRF.}}

\newglossaryentry{instantngp}{type=modeles, name={Instant-NGP}, sort={NeRF 1 Instant},
  parent=nerf,
  description={NeRF accéléré par une table de hachage multi-résolution et de petits MLP, avec une
  grille d'occupation \horscorpus. Version de départ de Tiny-DroNeRF : 527\,MB, 17,52\,GOp par
  pas, 32,59\,dB de PSNR \citep[Table~II]{arxiv-2603.01850-tiny-dronerf-tiny-neural-radiance}\pp.}}

\newglossaryentry{tiny-dronerf}{type=modeles, name={Tiny-DroNeRF}, sort={NeRF 2 Tiny},
  parent=instantngp,
  description={Instant-NGP réduit (1 image par pas, 160$\times$160, table de hachage $2^{13}$) et
  écrit en C pour le cluster du SoC GAP9 d'un nano-drone, en float16
  \citep[§IV]{arxiv-2603.01850-tiny-dronerf-tiny-neural-radiance}\pp. 0,284\,MB de paramètres,
  26,94\,dB de PSNR (5,7\,dB de moins), 73\,ms par pas d'entraînement amorti, 4,1\,s par rendu
  de 160$\times$160, mesurés au simulateur cycle-accurate GVSOC et non sur la puce
  \citep[§V-C]{arxiv-2603.01850-tiny-dronerf-tiny-neural-radiance}\pp ; 21,2\,dB avec des
  images de vol réel \citep[§V-D]{arxiv-2603.01850-tiny-dronerf-tiny-neural-radiance}\pp.}}

% --------------------------------------------------------------- Early-exit
\newglossaryentry{blankskip}{type=modeles, name={BlankSkip}, sort={BlankSkip},
  description={Sortie anticipée (\emph{early exit}) pour détecteur : une branche légère (0,078\,M
  de paramètres) accrochée à une couche intermédiaire d'un SSD-MobileNetV2 classe l'image
  \enquote{vide} ou \enquote{non vide} et court-circuite le détecteur dans le premier cas ;
  le seuil est réglable en vol \citep[§3]{arxiv-2603.28149-blankskip-early-exit-object-detection}\pp.
  Sur GAP8 (INT8, DORY) : 54,0 de mAP à 1,86\,FPS en moyenne contre 55,5 à 1,50\,FPS pour le
  modèle statique (+24\,\%) \citep[Table~2]{arxiv-2603.28149-blankskip-early-exit-object-detection}\pp ;
  2,07\,FPS sur vidéo d'exploration \citep[§4.2.1]{arxiv-2603.28149-blankskip-early-exit-object-detection}\pp.
  Puissance non mesurée.}}

% --------------------------------------------------------------- Opérateurs
\newglossaryentry{nms}{type=modeles, name={NMS}, sort={NMS},
  description={\emph{Non-Maximum Suppression} : suppression des boîtes redondantes d'un détecteur
  par seuil d'IoU \horscorpus. Dans le corpus c'est souvent le goulot : sur ARM, Nguyen et al. le
  mesurent à environ 3,9\,ms par image \citep[Table~12]{2024-nguyen-fpga-soc-implementation-yolov4-flying}, et
  \citet{arxiv-2607.13106-no-attention-no-problem-dpu}\pp concluent que le post-traitement sur
  l'ARM, non le DPU, limite le débit de bout en bout \citep[§IV-C]{arxiv-2607.13106-no-attention-no-problem-dpu}\pp.
  YOLO26 est annoncé sans NMS, mais celui-ci est conservé pour l'évaluation
  \citep[§III-B3]{arxiv-2607.13106-no-attention-no-problem-dpu}\pp ; Soft-NMS dans LCAM-YOLOX.}}

\newglossaryentry{attention}{type=modeles, name={Attention / Transformer}, sort={Attention},
  description={Mécanisme de pondération de positions par produit de requêtes et de clés,
  empilé en Transformer \horscorpus. Apparaît dans le corpus drone sous trois formes : têtes
  Transformer ou Swin de TPH-YOLOv5, Transformer de suivi (HiFT, TCTrack), attention de canal et
  spatiale (CBAM dans LCAM-YOLOX). Sur FPGA, elle pose problème : le module C2PSA de YOLO11 doit
  être approché par des convolutions 1$\times$1 et un hard-sigmoid pour passer sur le DPU
  \citep[§III-B]{arxiv-2607.13106-no-attention-no-problem-dpu}\pp.}}

\newglossaryentry{p2head}{type=modeles, name={Tête P2}, sort={Tete P2},
  description={Tête de détection supplémentaire sur la carte de stride 4 (160$\times$160 à
  640 px), dédiée aux objets minuscules. Recette récurrente des YOLO drone : TPH-YOLOv5,
  Drone-YOLO, Efficient YOLOv7-Drone, GCL-YOLO, DEAL-YOLO, QueryDet, DroneScan-YOLO. Gain
  important (43,7 vers 50,5 de mAP50 pour Drone-YOLO) mais coût élevé en activations :
  13,68 vers 10,89\,FPS pour TPH-YOLOv5 \citep[Table~5]{2023-zhao-tph-yolov5-boosting-object-detection}
  et latence TX2 multipliée par 4,1 pour Drone-YOLO nano
  \citep[Table~6]{2023-zhang-drone-yolo-efficient-neural-network}. Absente des YOLO portés sur FPGA
  du corpus ; SF-YOLOv5, lui, retire P5 sans ajouter P2.}}

\newglossaryentry{ghostconv}{type=modeles, name={GhostConv}, sort={GhostConv},
  description={Convolution \enquote{fantôme} (GhostNet) : une convolution réduite produit une
  partie des canaux, le reste étant généré par des opérations linéaires peu coûteuses
  \horscorpus. Réduit paramètres et FLOPs : backbone de GCL-YOLO, tête de YOLO-Drone, C3Ghost de
  GL-YOMO. Seule, elle fait pourtant chuter le mAP50 de GCL-YOLO de 32,7 à 28,8
  \citep[§4.5]{2023-cao-gcl-yolo-ghostconv-lightweight-yolo}, et elle est moins efficace sur DPU
  \deduit.}}
